\section{Conceptual and Legal Distinctions}

The tables below consolidate the principal distinctions employed across the argument. The first table presents a four-term analytical synthesis prepared for this essay to separate ontological dependence from moral duty, positive precept, and personal guilt. The second and third tables summarize classical distinctions from St.\ Thomas Aquinas's \emph{Summa Theologiae}.

\subsection*{Four-Term Distinction: Dependence, Duty, Precept, and Culpability}

\begin{fourstable}{Term}{Order}{What it establishes}{What it does not entail}
Dependence & Ontological fact & The creature continually receives being and every power of action from God alone. & Does not by itself specify a ceremonial rite or establish personal guilt. \\
Duty & Moral due (natural and revealed) & Rational creatures owe God acknowledgment, gratitude, and reverence freely in truth. & Does not authorize physical coercion in religious matters; binds in conscience, not by force. \\
Precept & Positive norm (divine or ecclesiastical) & Competent authority specifies concrete common forms, times, and sacraments for the common good. & Does not make mutable human discipline eternal dogma, nor apply outside its declared jurisdiction. \\
Culpability & Personal responsibility & Imputes formal guilt only where objective matter, sufficient knowledge, and deliberate freedom concur. & Does not follow automatically from material non-observance under illness, ignorance, or grave necessity. \\
\end{fourstable}

\subsection*{Acts of Religion and Levels of Moral Action}

\begin{threestable}{Category}{Distinction in Aquinas}{Significance for worship}
Elicited interior acts & Devotion and prayer (\emph{ST} II-II, qq.~82--83) & Belong essentially and primarily to religion; express the readiness of will and raising of the mind to God. \\
Elicited exterior acts & Adoration and sacrifice (\emph{ST} II-II, qq.~84--85) & Secondary signs ordered to interior reverence; engage the body and sensible creation in confessing subjection to God. \\
Commanded acts & Works of mercy, temperance, and justice directed to God's honor (\emph{ST} II-II, q.~81, a.~1, ad~1; Jas~1:27) & Elicited by other virtues (such as mercy toward the poor) yet directed by religion to God as end. \\
Objective conformity & Fulfilling the substance of a precept (\emph{ST} I-II, q.~100, a.~10) & A genuine good in its kind even before warm affection matures; avoids transgressing the specific precept. \\
Acquired moral virtue & Prompt, steady, well-ordered action (\emph{ST} I-II, q.~107, a.~4) & Formed over time through repeated due acts, educating appetite and body toward willing ease. \\
Supernatural merit & Free cooperation in sanctifying grace animated by charity (\emph{ST} I-II, q.~114, aa.~1--2, 5) & Rests entirely on God's gratuitous ordination in Christ; never earned prior to grace and never puts God in debt. \\
\end{threestable}

\clearpage

\subsection*{Kinds of Law Ordering Worship}

\begin{fourstable}{Kind of law}{Source in Aquinas}{Scope and basis}{Role in Christian worship}
Eternal law & \emph{ST} I-II, q.~91, a.~1 & The eternal plan of Divine Wisdom governing the whole universe toward its end. & Ultimate fountain of all right order and all valid human and divine norms. \\
Natural law & \emph{ST} I-II, q.~91, a.~2; q.~94, a.~2 & The rational creature's participation in eternal law through practical reason. & Recognizes the duty to honor God, give time to worship, and offer outward signs of subjection (\emph{ST} II-II, q.~85, a.~1). \\
Divine law (Old and New) & \emph{ST} I-II, q.~91, a.~4; qq.~106--108 & Revealed by God; the New Law is chiefly the Spirit's grace. & Reveals Christ's unique oblation, institutes the seven sacraments, and governs interior rectitude as a law of liberty. \\
Human and Church law & \emph{ST} I-II, q.~91, a.~3; q.~108, a.~2 & Determinate norms set by competent authority for the common good. & Fixes shared liturgical times and common order (e.g., universal Latin Sunday discipline in CIC 1983 cc.~1245--1248). \\
\end{fourstable}
